\section{Main}\label{sec:introduction}

% TODO: Figure 1 should illustrate: (a) disease-associated vs disease-relevant distinction,
% (b) PIGEAN framework overview (graphical model), (c) example gene positioned along
% direct and indirect support axes, (d) downstream applications

% P1: Broad - gene-trait mapping matters, genetics is powerful but incomplete
Most human traits are influenced by many genes, and most human genes influence many traits. Mapping this complex web of gene--trait relationships is essential for uncovering disease biology \cite{cohen_sequence_2006}, understanding clinical heterogeneity \cite{udler_type_2018}, accurately diagnosing patients \cite{yang_clinical_2013}, and designing effective therapies \cite{plenge_validating_2013}. Among the most widely used and powerful approaches to identify gene-trait relationships are genetic association studies. For common traits, genome-wide association studies (GWAS) have identified tens of thousands of trait-associated loci across thousands of phenotypes \cite{abdellaoui_15_2023}, while rare disease studies have linked thousands of Mendelian conditions to specific causal genes \cite{chung_meta-analysis_2023}. A genetic association provides a gene with ``genetic support'' for a human trait---establishing the human-relevance and causal effect of its perturbation.

% P2: The gap - associated vs relevant, most relevant genes escape detection
The number of trait-associated genes identified by genetic association studies is, however, but a small fraction of the total number of trait-relevant genes---that is, those that can be perturbed in humans to yield an effect on a trait of interest \cite{finan_druggable_2017}. Many trait-relevant genes escape detection due to limited statistical power \cite{wang_statistical_2019}, constraining evolutionary pressure on genetic variation \cite{gazal_linkage_2017}, or the difficulty of assigning noncoding GWAS signals to their effector genes \cite{schipper_demystifying_2022,costanzo_realizing_2025}. Yet the full set of trait-relevant genes is of most interest for drug development and experimental study \cite{jhamb_pathway_2019,macnamara_network_2020}. In part due to this incompleteness, most translational research does not use human genetics to prioritize genes \cite{wehling_assessing_2009,dornbos_evaluating_2022}, even with the increasing recognition that drug targets with genetic support are substantially more likely to succeed clinically \cite{plenge_validating_2013,minikel_refining_2024}.

% P3: Non-genetic approaches - widely used but biased
By contrast, non-human-genetic approaches are more widely used to identify trait-relevant genes but come with critical tradeoffs. Experimental studies in animal \cite{seok_genomic_2013,widjaja_inhibiting_2019} and cell-based \cite{ben-david_genetic_2018,canon_clinical_2019} systems can establish causal links between genes and phenotypes, but their relevance to humans is often uncertain. Human observational studies, such as gene expression profiling \cite{wu_integrating_2022,rodriguez_machine_2021} or epidemiological associations \cite{barter_cholesteryl_2011,bhatt_cardiovascular_2019}, provide human-relevant context but lack causal directionality. The literature and biomedical databases are rife with gene-trait relationships derived from these or other non-genetic studies, and they are widely used as a starting point for drug discovery or experimental study \cite{swinney_how_2011,dornbos_evaluating_2022}. Recent studies demonstrate how these sources can be rapidly and accurately queried with large language models \cite{toufiq_harnessing_2023,hu_evaluation_2025}. However, it is well documented that prioritizing genes using these evidence sources is prone to bias, incompleteness, or circularity \cite{gillis_assessing_2013,haynes_gene_2018,stoeger_large_2018,tan_caution_2025}.

% P4: Existing integration methods are heuristic
Ideally, a map of gene-trait relationships would combine the complementary benefits of human genetic and non-human-genetic evidence sources. It is commonplace to evaluate genes individually through manual expert synthesis of genetic data and the published literature \cite{mountjoy_open_2022,karjalainen_genome-wide_2024}, but a quantitative map across the genome and phenome requires scalable statistical methods. Such methods have been developed both in the context of GWAS effector gene identification, where gene functional annotations enriched for genetic associations are used to prioritize genes nearby associated SNPs \cite{pers_biological_2015,weeks_leveraging_2023}, and in the context of rare disease diagnostics, where literature mining \cite{birgmeier_amelie_2020} and interaction networks \cite{smedley_next-generation_2015} are used to prioritize genes that carry variants observed in a patient. These approaches, however, rely on heuristics (such as harmonic averaging \cite{buniello_open_2025} or regressions using noisy or incomplete training data \cite{weeks_leveraging_2023,schipper_prioritizing_2025}) and often lack a biological model, which limits their portability across traits or application to settings other than that for which they were originally designed. As a result, the currently charted map of trait-gene relationships across rare and common diseases remains sparse, inaccurate, and inconsistent across traits.

% P5: Our contribution - PIGEAN framework, calibration, and resource
Here, we combine human genetic association data for rare and common traits and non-human-genetic data encoded in {{NUM_TOTAL_GENE_ANNOTATIONS:,}} gene annotations to construct a probabilistic evidence map relating 18,125 genes to {{NUM_TOTAL_TRAITS_MOUSE_MSIGDB:,}} traits ({{NUM_COMMON_TRAITS_MOUSE_MSIGDB:,}} common, {{NUM_RARE_TRAITS_MOUSE_MSIGDB:,}} rare). Each gene-trait relationship is a composite of two statistically calibrated values: a traditional measure of \textit{direct genetic support} (the strength of gene-trait association from human genetic data) and a new measure of \textit{indirect genetic support} (the extent to which the gene shares functional annotations with other trait-associated genes). Indirect genetic support is measured on the same scale as direct genetic support, estimates a measure of gene-trait relevance distinct from direct gene-trait association, and enables us to use unbiased human genetic associations to calibrate noisy and heterogeneous non-human-genetic data sources and learn which are trustworthy for each trait. We validate indirect support on three independent axes: it predicts held-out genetic associations in leave-one-chromosome-out and temporal analyses, prioritizes drug targets with higher clinical success rates than existing resources, and enriches for causal genes in rare disease patients. This framework provides a probabilistic map of gene-trait relevance that is orders of magnitude denser than manually curated knowledge graphs, and a principled replacement for the ad hoc gene prioritization approaches currently used across the field.

% ANALYSIS IDEA: Consider adding a "headline number" for how many additional drug targets
% indirect support surfaces compared to direct-only approaches at equivalent success rates.
% This would strengthen the translational impact claim in the contribution paragraph.
